%% toc: P2 22 22.1,22.2,22.3
\section{Proofs for the sort conventions and the one-sorted reductions}
\label{sec:annex-p02-w02-p02-038-many-sorted-reduction}

Notation is that of \cref{sec:p02-w02-p02-038-many-sorted-reduction}.  Sorted
variables are identified with their images in the one-sorted variable
supply.  For a one-sorted assignment \(\sigma\), a variable \(v\) and an
element \(d\), \(\sigma[v\mapsto d]\) is the one-sorted update.  The guard
translation commutes with \(\bottom,\neg,\land,\lor,\to\), leaves atoms
unchanged, and guards quantifiers as in
\cref{def:p02-many-sorted-translation}.

\begin{annexitem}{annex:p02-w02-p02-038-coincidence}{thm:p02-w02-p02-038-coincidence}{Coincidence for contextual assignments}
\begin{proof}
\emph{Terms}, by induction on \(t\), for all \(g,g'\) agreeing on every
variable of \(t\).  If \(t\) is a variable \(v\), then
\(\llbracket v\rrbracket^{\mathcal M}_{g}=g(v)=g'(v)=
\llbracket v\rrbracket^{\mathcal M}_{g'}\).  If \(t\) is a constant \(c\),
both values are \(c^{\mathcal M}\).  If \(t=f(t_1,\ldots,t_n)\), the
variables of each \(t_i\) are variables of \(t\), so the induction hypothesis
gives equal argument values and therefore equal values of \(f^{\mathcal M}\).

\emph{Formulas}, by induction on \(\varphi\), for all \(g,g'\) whose domains
contain \(\FV(\varphi)\) and which agree on \(\FV(\varphi)\).  For
\(P(t_1,\ldots,t_n)\) and for \(t=u\), the free variables are the variables
of the terms, so the term case gives equal values and the two atomic clauses
give the same verdict.  \(\bottom\) is false under both.  For \(\neg\psi\)
and for \(\psi\land\chi\), \(\psi\lor\chi\), \(\psi\to\chi\), the free
variables of each component are among those of \(\varphi\); the induction
hypothesis gives the same verdict on each component, and the clauses combine
the verdicts in the same way.  Let \(\varphi=\forall v\,\psi\) or
\(\exists v\,\psi\) with \(v\) of sort \(\mathsf s\), and let
\(b\in M_{\mathsf s}\).  The assignments \(g[v\mapsto b]\) and
\(g'[v\mapsto b]\) both send \(v\) to \(b\), and agree on
\(\FV(\psi)\setminus\{v\}\subseteq\FV(\varphi)\); so they agree on
\(\FV(\psi)\), and the induction hypothesis gives
\(\mathcal M\entails\psi[g[v\mapsto b]]\) iff
\(\mathcal M\entails\psi[g'[v\mapsto b]]\).  As this holds for every
\(b\in M_{\mathsf s}\), the condition ``for every \(b\)'' has the same verdict
for \(g\) and \(g'\), and so has the condition ``for some \(b\)''.

For the final clause, \(g'\) and \(g'{\restriction}\Delta\) agree on
\(\Delta\supseteq\FV(\varphi)\).
\end{proof}
\end{annexitem}

\begin{annexitem}{annex:p02-w02-p02-038-guard-nonempty}{thm:p02-w02-p02-038-guard-nonempty}{The assignment guard implies the nonemptiness family}
\begin{proof}
Let \(\mathsf s\) be a sort.  Its designated variable \(x^{\mathsf s}\) is in
\(\mathrm{Var}_{\mathrm{MS}}\), so \(Q_{\mathsf s}(x^{\mathsf s})\) is a
member of \(\mathsf{VarGuard}_{\mathrm{MS}}\) and
\(\sigma(x^{\mathsf s})\in Q_{\mathsf s}^{\mathcal A}\).  This element
witnesses \(\mathcal A\entails\exists x\,Q_{\mathsf s}(x)\).

Since
\(\mathsf{Ax}^{\mathrm E}_{\mathrm{sort}}\subseteq\mathsf{Ax}_{\mathrm{sort}}\),
a consequence of the smaller premise set is a consequence of the larger one.
Conversely, suppose \(\mathsf{Ax}_{\mathrm{sort}}\cup
\mathsf{VarGuard}_{\mathrm{MS}}\cup X\entails_{\mathrm{FOL}}\chi\), and let
\(\mathcal A,\sigma\) satisfy every member of
\(\mathsf{Ax}^{\mathrm E}_{\mathrm{sort}}\cup\mathsf{VarGuard}_{\mathrm{MS}}
\cup X\).  By the first paragraph, \(\mathcal A\) satisfies every sentence of
\((\mathrm{Ne})\), and the remaining members of
\(\mathsf{Ax}_{\mathrm{sort}}\) are in
\(\mathsf{Ax}^{\mathrm E}_{\mathrm{sort}}\).  So \(\mathcal A,\sigma\)
satisfies the larger premise set and therefore \(\chi\).
\end{proof}
\end{annexitem}

\begin{annexitem}{annex:p02-w02-p02-038-coverage-not-elementary}{thm:p02-w02-p02-038-coverage-not-elementary}{Coverage is not first-order with infinitely many sorts}
\begin{proof}
First, a covering model.  Let \(\mathcal A_0\) have domain
\(\mathsf{Sort}\) (nonempty by \cref{def:p02-many-sorted-signature}), with
\(Q_{\mathsf s}^{\mathcal A_0}\defeq\{\mathsf s\}\),
\(c^{\mathcal A_0}\defeq\mathsf s\) for \(c:\longrightarrow\mathsf s\),
\(f^{\mathcal A_0}(d_1,\ldots,d_n)\defeq\mathsf s\) for every input when
\(f\) has codomain \(\mathsf s\), and every relation symbol empty.  Then
\((\mathrm{Ne})\) holds since \(\mathsf s\in Q_{\mathsf s}^{\mathcal A_0}\);
\((\mathrm{Ct})\) holds since \(c^{\mathcal A_0}=\mathsf s\); every value of
\(f^{\mathcal A_0}\) is \(\mathsf s\in Q_{\mathsf s}^{\mathcal A_0}\), so
\((\mathrm{Cl})\) holds; and \(\{\mathsf s\}\cap\{\mathsf t\}=\varnothing\)
for \(\mathsf s\ne\mathsf t\), so \((\mathrm{Dj})\) holds.  Every element
\(\mathsf s\) of the domain lies in \(Q_{\mathsf s}^{\mathcal A_0}\).

Suppose, for a contradiction, that \(T\) is as described.  Then
\(\mathcal A_0\entails T\).  Let \(e\) be a new constant symbol and put
\[
 \Sigma\defeq T\cup\mathsf{Ax}_{\mathrm{sort}}(\Lang_{\mathrm{MS}})\cup
 \{\neg Q_{\mathsf s}(e):\mathsf s\in\mathsf{Sort}\}.
\]
A finite \(\Sigma_0\subseteq\Sigma\) contains \(\neg Q_{\mathsf s}(e)\) only
for \(\mathsf s\) in a finite set \(S_0\).  As \(\mathsf{Sort}\) is
infinite, there is \(\mathsf t\in\mathsf{Sort}\setminus S_0\).  Expand
\(\mathcal A_0\) by \(e\mapsto\mathsf t\).  The expansion satisfies \(T\) and
\(\mathsf{Ax}_{\mathrm{sort}}\), which do not mention \(e\), and
\(\neg Q_{\mathsf s}(e)\) for \(\mathsf s\in S_0\), because
\(\mathsf t\notin\{\mathsf s\}\).  So \(\Sigma_0\) has a model.  By
\cref{thm:fol-compactness}, clause~1 (which uses the Axiom of Choice),
\(\Sigma\) has a model \(\mathcal B\).  By \cref{lem:reduct-invariance}, the
\(\Lang_{\mathrm{MS}}^{\flat}\)-reduct of \(\mathcal B\) satisfies \(T\) and
\(\mathsf{Ax}_{\mathrm{sort}}\), so by the property of \(T\) the element
\(e^{\mathcal B}\) lies in \(Q_{\mathsf t}^{\mathcal B}\) for some sort
\(\mathsf t\).  This contradicts \(\mathcal B\entails\neg Q_{\mathsf t}(e)\).
\end{proof}
\end{annexitem}

\begin{annexitem}{annex:p02-w02-p02-038-restriction}{thm:p02-w02-p02-038-restriction}{Restriction lemma}
\begin{proof}
\emph{Clause 1.}  For \(c:\longrightarrow\mathsf s\), \((\mathrm{Ct})\) gives
\(c^{\mathcal A}\in Q_{\mathsf s}^{\mathcal A}=M_{\mathsf s}\).  For
\(f:\mathsf s_1\times\cdots\times\mathsf s_n\longrightarrow\mathsf s\) and
\(b_i\in M_{\mathsf s_i}=Q_{\mathsf s_i}^{\mathcal A}\), the instance of
\((\mathrm{Cl})\) for \(f\) at \(b_1,\ldots,b_n\) has a true antecedent, so
\(f^{\mathcal A}(b_1,\ldots,b_n)\in Q_{\mathsf s}^{\mathcal A}=M_{\mathsf s}\);
hence the restriction of \(f^{\mathcal A}\) maps
\(M_{\mathsf s_1}\times\cdots\times M_{\mathsf s_n}\) into \(M_{\mathsf s}\).
By definition \(P^{\mathcal A^{\downarrow}}\) is a subset of the right
product.  So \(\mathcal A^{\downarrow}\) is a sorted structure.  If
\(\mathcal A\entails(\mathrm{Dj})\), no element lies in
\(Q_{\mathsf s}^{\mathcal A}\) and \(Q_{\mathsf t}^{\mathcal A}\) for
\(\mathsf s\ne\mathsf t\), so the carriers are pairwise disjoint.  If
\(\mathcal A\entails(\mathrm{Ne})\), each
\(M_{\mathsf s}=Q_{\mathsf s}^{\mathcal A}\) has an element.  Both together
give strictness.  Finally, \(\mathsf{Guard}(\Delta)\) gives
\(\sigma(v)\in Q_{\mathsf s}^{\mathcal A}=M_{\mathsf s}\) for every
\(v\in\Delta\) of sort \(\mathsf s\), so \(\sigma{\restriction}\Delta\) is a
\(\Delta\)-assignment.

\emph{Clause 2}, by induction on \(t\); write \(g\defeq\sigma{\restriction}\Delta\).
If \(t=v\in\Delta\) has sort \(\mathsf s\), both values are \(\sigma(v)\),
which is in \(Q_{\mathsf s}^{\mathcal A}\) by the guard.  If \(t=c\), both
values are \(c^{\mathcal A}\), in \(Q_{\mathsf s}^{\mathcal A}\) by
\((\mathrm{Ct})\).  If \(t=f(t_1,\ldots,t_n)\), the induction hypothesis
gives
\(\llbracket t_i\rrbracket^{\mathcal A}_{\sigma}=
\llbracket t_i\rrbracket^{\mathcal A^{\downarrow}}_{g}\eqqcolon b_i\in
M_{\mathsf s_i}\).  Then
\(\llbracket t\rrbracket^{\mathcal A}_{\sigma}=f^{\mathcal A}(b_1,\ldots,b_n)
=f^{\mathcal A^{\downarrow}}(b_1,\ldots,b_n)=
\llbracket t\rrbracket^{\mathcal A^{\downarrow}}_{g}\), the middle equation
because the tuple lies in the restricted domain, and this value lies in
\(M_{\mathsf s}=Q_{\mathsf s}^{\mathcal A}\) by clause~1.

\emph{Clause 3}, by induction on \(\varphi\), simultaneously for every
context \(\Delta\supseteq\FV(\varphi)\) and every \(\sigma\) satisfying
\(\mathsf{Guard}(\Delta)\); again \(g\defeq\sigma{\restriction}\Delta\).
\begin{itemize}
\item \(P(t_1,\ldots,t_n)\).  By clause~2 the tuple
      \((\llbracket t_i\rrbracket^{\mathcal A}_{\sigma})_i\) equals
      \((\llbracket t_i\rrbracket^{\mathcal A^{\downarrow}}_{g})_i\) and lies
      in \(M_{\mathsf s_1}\times\cdots\times M_{\mathsf s_n}\).  A tuple of
      that product is in \(P^{\mathcal A}\) iff it is in
      \(P^{\mathcal A}\cap(M_{\mathsf s_1}\times\cdots\times M_{\mathsf s_n})
      =P^{\mathcal A^{\downarrow}}\).  So
      \(\mathcal A,\sigma\entails P(t_1,\ldots,t_n)\) iff
      \(\mathcal A^{\downarrow}\entails P(t_1,\ldots,t_n)[g]\).
\item \(t=u\), of the same or of different sorts.  By clause~2 the two
      one-sorted values are the two sorted values, so they are equal in
      \(\mathcal A\) iff they are equal in \(\mathcal A^{\downarrow}\).
\item \(\bottom\).  Its translation is \(\bottom\), false on both sides.
\item \(\neg\psi\).  \(\FV(\psi)=\FV(\varphi)\subseteq\Delta\).
      \(\mathcal A,\sigma\entails\neg\operatorname{tr}_{\mathrm{MS}}(\psi)\)
      iff \(\mathcal A,\sigma\not\entails\operatorname{tr}_{\mathrm{MS}}(\psi)\)
      iff, by the induction hypothesis,
      \(\mathcal A^{\downarrow}\not\entails\psi[g]\), iff
      \(\mathcal A^{\downarrow}\entails\neg\psi[g]\).
\item \(\psi\land\chi\), \(\psi\lor\chi\), \(\psi\to\chi\).  Both
      components have free variables in \(\Delta\).  The induction hypothesis
      gives the same verdict for \(\operatorname{tr}_{\mathrm{MS}}(\psi)\)
      under \(\sigma\) as for \(\psi\) under \(g\), and likewise for
      \(\chi\); the one-sorted and the sorted clauses for each connective
      combine the two verdicts by the same truth function.
\item \(\forall v\,\psi\) with \(v\) of sort \(\mathsf s\).  Its translation
      \(\forall v\,(Q_{\mathsf s}(v)\to\operatorname{tr}_{\mathrm{MS}}(\psi))\)
      holds under \(\sigma\) iff, for every \(d\) in the domain, either
      \(d\notin Q_{\mathsf s}^{\mathcal A}\) or
      \(\mathcal A,\sigma[v\mapsto d]\entails\operatorname{tr}_{\mathrm{MS}}(\psi)\);
      that is, iff the second alternative holds for every
      \(d\in Q_{\mathsf s}^{\mathcal A}=M_{\mathsf s}\).  For such \(d\), the
      assignment \(\sigma[v\mapsto d]\) satisfies
      \(\mathsf{Guard}(\Delta\cup\{v\})\): at \(v\) its value is
      \(d\in Q_{\mathsf s}^{\mathcal A}\), and at every other member of
      \(\Delta\) it agrees with \(\sigma\).  Also
      \(\FV(\psi)\subseteq\Delta\cup\{v\}\) and
      \((\sigma[v\mapsto d]){\restriction}(\Delta\cup\{v\})=g[v\mapsto d]\).
      The induction hypothesis for \(\psi\) at the context
      \(\Delta\cup\{v\}\) therefore gives
      \(\mathcal A,\sigma[v\mapsto d]\entails\operatorname{tr}_{\mathrm{MS}}(\psi)\)
      iff \(\mathcal A^{\downarrow}\entails\psi[g[v\mapsto d]]\).  So the
      translation holds iff \(\mathcal A^{\downarrow}\entails\psi[g[v\mapsto d]]\)
      for every \(d\in M_{\mathsf s}\), which is the clause for
      \(\mathcal A^{\downarrow}\entails\forall v\,\psi[g]\).
\item \(\exists v\,\psi\) with \(v\) of sort \(\mathsf s\).  Its translation
      \(\exists v\,(Q_{\mathsf s}(v)\land\operatorname{tr}_{\mathrm{MS}}(\psi))\)
      holds under \(\sigma\) iff some \(d\in Q_{\mathsf s}^{\mathcal A}=
      M_{\mathsf s}\) has
      \(\mathcal A,\sigma[v\mapsto d]\entails\operatorname{tr}_{\mathrm{MS}}(\psi)\).
      For each such \(d\) the facts established in the previous case about
      \(\sigma[v\mapsto d]\) hold verbatim, so the induction hypothesis at
      \(\Delta\cup\{v\}\) turns this into: some \(d\in M_{\mathsf s}\) has
      \(\mathcal A^{\downarrow}\entails\psi[g[v\mapsto d]]\), which is the
      clause for \(\mathcal A^{\downarrow}\entails\exists v\,\psi[g]\).
\end{itemize}
\end{proof}
\end{annexitem}

\begin{annexitem}{annex:p02-w02-p02-038-transport}{thm:p02-w02-p02-038-transport}{Transport lemma}
\begin{proof}
\emph{Well-definedness.}  If \(d_i=h_{\mathsf s_i}(b_i)=h_{\mathsf s_i}(b_i')\)
then \(b_i=b_i'\) by injectivity, so the first case of \(f^{\mathcal M^{h}}\)
names one value; every other tuple takes the fill value, so
\(f^{\mathcal M^{h}}\) is a total function \(D^n\to D\).  Constants take
values in \(D\), relations are subsets of \(D^n\), and \(D\ne\varnothing\);
so \(\mathcal M^{h}\) is a one-sorted \(\Lang_{\mathrm{MS}}^{\flat}\)-structure.

\emph{Clause 1}, by induction on \(t\).  For \(v\in\Delta\) of sort
\(\mathsf s\), \(\sigma(v)=h_{\mathsf s}(g(v))\) because \(\sigma\)
transports \(g\).  For \(c:\longrightarrow\mathsf s\),
\(c^{\mathcal M^{h}}=h_{\mathsf s}(c^{\mathcal M})\) by definition.  For
\(t=f(t_1,\ldots,t_n)\), put
\(b_i\defeq\llbracket t_i\rrbracket^{\mathcal M}_{g}\in M_{\mathsf s_i}\).
The induction hypothesis gives
\(\llbracket t_i\rrbracket^{\mathcal M^{h}}_{\sigma}=h_{\mathsf s_i}(b_i)\),
so the argument tuple falls under the first case, and
\(\llbracket t\rrbracket^{\mathcal M^{h}}_{\sigma}=
h_{\mathsf s}(f^{\mathcal M}(b_1,\ldots,b_n))=
h_{\mathsf s}(\llbracket t\rrbracket^{\mathcal M}_{g})\).  Fill values are
never reached by a term evaluated under \(\sigma\).

\emph{Clause 2}, by induction on \(\varphi\), simultaneously for every
context \(\Delta\supseteq\FV(\varphi)\), every \(\Delta\)-assignment \(g\) and
every \(\sigma\) transporting \(g\).
\begin{itemize}
\item \(P(t_1,\ldots,t_n)\).  Put
      \(b_i\defeq\llbracket t_i\rrbracket^{\mathcal M}_{g}\); by clause~1 the
      one-sorted argument tuple is \((h_{\mathsf s_i}(b_i))_i\).  If
      \((b_i)_i\in P^{\mathcal M}\), this tuple is in \(P^{\mathcal M^{h}}\)
      by definition.  If the tuple is in \(P^{\mathcal M^{h}}\), there is
      \((b_i')_i\in P^{\mathcal M}\) with
      \(h_{\mathsf s_i}(b_i')=h_{\mathsf s_i}(b_i)\) for each \(i\); by
      injectivity \(b_i'=b_i\), so \((b_i)_i\in P^{\mathcal M}\).
\item \(t=u\) with both terms of sort \(\mathsf s\).  With
      \(a\defeq\llbracket t\rrbracket^{\mathcal M}_{g}\) and
      \(b\defeq\llbracket u\rrbracket^{\mathcal M}_{g}\), clause~1 makes the
      one-sorted values \(h_{\mathsf s}(a)\) and \(h_{\mathsf s}(b)\), which
      are equal iff \(a=b\) by injectivity.
\item \(t=u\) with \(t\) of sort \(\mathsf s\) and \(u\) of sort
      \(\mathsf t\ne\mathsf s\) (cross-sort language, \(h\) coherent).  The
      one-sorted values are \(h_{\mathsf s}(a)\) and \(h_{\mathsf t}(b)\),
      equal iff \(a=b\) by coherence.
\item \(\bottom\), \(\neg\psi\), and \(\psi\land\chi\), \(\psi\lor\chi\),
      \(\psi\to\chi\).  The translation is \(\bottom\), false on both sides,
      or the same connective applied to the translated components.  The
      components have free variables in \(\Delta\), belong to the same
      language as \(\varphi\), and the induction hypothesis gives each the
      same verdict on both sides; the connective clauses combine verdicts by
      the same truth function.
\item \(\forall v\,\psi\) with \(v\) of sort \(\mathsf s\).  The translation
      holds under \(\sigma\) iff
      \(\mathcal M^{h},\sigma[v\mapsto d]\entails\operatorname{tr}_{\mathrm{MS}}(\psi)\)
      for every \(d\in Q_{\mathsf s}^{\mathcal M^{h}}=h_{\mathsf s}[M_{\mathsf s}]\),
      that is, iff
      \(\mathcal M^{h},\sigma[v\mapsto h_{\mathsf s}(b)]\entails
      \operatorname{tr}_{\mathrm{MS}}(\psi)\) for every \(b\in M_{\mathsf s}\).
      The assignment \(\sigma[v\mapsto h_{\mathsf s}(b)]\) transports
      \(g[v\mapsto b]\): at \(v\) both sides are \(h_{\mathsf s}(b)\), and at
      \(w\in\Delta\setminus\{v\}\) of sort \(\mathsf t\) the value is
      \(\sigma(w)=h_{\mathsf t}(g(w))\).  Since
      \(\FV(\psi)\subseteq\Delta\cup\{v\}\), the induction hypothesis at that
      context gives
      \(\mathcal M^{h},\sigma[v\mapsto h_{\mathsf s}(b)]\entails
      \operatorname{tr}_{\mathrm{MS}}(\psi)\) iff
      \(\mathcal M\entails\psi[g[v\mapsto b]]\).  So the translation holds iff
      \(\mathcal M\entails\psi[g[v\mapsto b]]\) for every
      \(b\in M_{\mathsf s}\), the clause for
      \(\mathcal M\entails\forall v\,\psi[g]\).
\item \(\exists v\,\psi\) with \(v\) of sort \(\mathsf s\).  The translation
      holds under \(\sigma\) iff some \(d\in h_{\mathsf s}[M_{\mathsf s}]\),
      equivalently some \(h_{\mathsf s}(b)\) with \(b\in M_{\mathsf s}\),
      has \(\mathcal M^{h},\sigma[v\mapsto h_{\mathsf s}(b)]\entails
      \operatorname{tr}_{\mathrm{MS}}(\psi)\).  The transport fact of the
      previous case holds for each such \(b\), and the induction hypothesis
      turns the condition into: some \(b\in M_{\mathsf s}\) has
      \(\mathcal M\entails\psi[g[v\mapsto b]]\), the clause for
      \(\mathcal M\entails\exists v\,\psi[g]\).
\end{itemize}

\emph{Clause 3.}  \((\mathrm{Ct})\): \(c^{\mathcal M^{h}}=
h_{\mathsf s}(c^{\mathcal M})\in h_{\mathsf s}[M_{\mathsf s}]=
Q_{\mathsf s}^{\mathcal M^{h}}\).  \((\mathrm{Cl})\): let
\(d_i\in Q_{\mathsf s_i}^{\mathcal M^{h}}\) for every \(i\).  Then
\(d_i=h_{\mathsf s_i}(b_i)\) with \(b_i\in M_{\mathsf s_i}\), the first case
applies, and \(f^{\mathcal M^{h}}(d_1,\ldots,d_n)=
h_{\mathsf s}(f^{\mathcal M}(b_1,\ldots,b_n))\in Q_{\mathsf s}^{\mathcal M^{h}}\).
A fill value is produced only at a tuple with some
\(d_i\notin Q_{\mathsf s_i}^{\mathcal M^{h}}\), where the antecedent of the
closure axiom is false; so the axiom holds whatever the fill functions are.
\((\mathrm{Ne})\): \(h_{\mathsf s}[M_{\mathsf s}]\ne\varnothing\) iff
\(M_{\mathsf s}\ne\varnothing\).  \((\mathrm{Dj})\): its instance for
\(\mathsf s\ne\mathsf t\) holds iff no element lies in
\(h_{\mathsf s}[M_{\mathsf s}]\cap h_{\mathsf t}[M_{\mathsf t}]\).
\(\mathsf{Cov}\): it holds iff every element of \(D\) lies in some
\(h_{\mathsf s}[M_{\mathsf s}]\).

\emph{Clause 4.}  For \(v\in\Delta\) of sort \(\mathsf s\),
\(\sigma(v)=h_{\mathsf s}(g(v))\in Q_{\mathsf s}^{\mathcal M^{h}}\).

\emph{The three instances.}  (i) Tagged with junk point.  Each
\(h_{\mathsf s}(a)=(\mathsf s,a)\) is injective.  If
\((\mathsf s,a)=(\mathsf t,b)\) then \(\mathsf s=\mathsf t\), so the images
are pairwise disjoint.  By Foundation \(\mathsf{Sort}\notin\mathsf{Sort}\), so
\(\star=(\mathsf{Sort},\mathsf{Sort})\) differs from every \((\mathsf s,a)\)
with \(\mathsf s\in\mathsf{Sort}\) and lies in no
\(Q_{\mathsf s}^{\mathcal M^{\sharp}}\); in particular \(D\ne\varnothing\),
the closure axioms are not violated by the junk values (clause~3), and
\(\mathsf{Cov}\) fails at \(\star\).  The assignment sending each
\(v\in\Delta\) of sort \(\mathsf s\) to \((\mathsf s,g(v))\) and every other
one-sorted variable to \(\star\) transports \(g\).  If the carriers are
pairwise disjoint, the family is coherent: for \(\mathsf s\ne\mathsf t\),
\(a\in M_{\mathsf s}\), \(b\in M_{\mathsf t}\), both \(a=b\) and
\((\mathsf s,a)=(\mathsf t,b)\) are false, and for \(\mathsf s=\mathsf t\)
coherence is injectivity.  If some \(a\) lies in \(M_{\mathsf s}\cap
M_{\mathsf t}\) with \(\mathsf s\ne\mathsf t\), then \(a=a\) while
\((\mathsf s,a)\ne(\mathsf t,a)\), so the family is not coherent.

(ii) Tagged with defaults, for strict \(\mathcal M\) and total \(s\).  The
same \(h\) is injective with pairwise disjoint images.  \(D\) is nonempty
because \(\mathsf{Sort}\) and every carrier are nonempty, and
\(j_f\) takes the value \((\mathsf s,s(x^{\mathsf s}))\in D\) for \(f\) of
codomain \(\mathsf s\), since \(s(x^{\mathsf s})\in M_{\mathsf s}\).  Every
element of \(D\) is a tagged pair, so the images cover \(D\).  The
assignment sending each sorted \(v\) of sort \(\mathsf s\) to
\((\mathsf s,s(v))\) and every other one-sorted variable to the fixed
element transports \(s\).

(iii) Plain union with junk point.  Inclusions are injective, and
\(h_{\mathsf s}(a)=h_{\mathsf t}(b)\) means \(a=b\), so the family is
coherent.  By Foundation \(U\notin U\), so the junk point \(U\) is not in any
carrier and lies in no \(Q_{\mathsf s}^{\mathcal M^{\cup}}\); \(D\ne\varnothing\)
and \(\mathsf{Cov}\) fails at \(U\).  The images are the carriers
themselves, so they are pairwise disjoint iff the carriers are.  The
assignment sending \(v\in\Delta\) to \(g(v)\) and every other one-sorted
variable to \(U\) transports \(g\).
\end{proof}
\end{annexitem}

\begin{annexitem}{annex:p02-w02-p02-038-strict-correspondence}{thm:p02-w02-p02-038-strict-correspondence}{Two-way satisfaction correspondence for the strict reduction}
\begin{proof}
\emph{Clause 1.}  \(\mathcal A\) satisfies all four families, and
\[
 \mathsf{VarGuard}_{\mathrm{MS}}=\mathsf{Guard}(\mathrm{Var}_{\mathrm{MS}}).
\]
By \cref{thm:p02-w02-p02-038-restriction}, clause~1, \(\mathcal A^{\downarrow}\)
is strict and \(\sigma{\restriction}\mathrm{Var}_{\mathrm{MS}}\) is a total
assignment into it.  A well-sorted formula is a formula of the cross-sort
language with \(\FV(\varphi)\subseteq\mathrm{Var}_{\mathrm{MS}}\), so clause~3
of that lemma at \(\Delta=\mathrm{Var}_{\mathrm{MS}}\) gives the equivalence;
for a strict structure and a total assignment, the contextual relation is
the relation of \cref{def:p02-many-sorted-structure}
(\cref{def:p02-w02-p02-038-satisfaction}).

\emph{Clause 2.}  \(\mathcal M^{\sharp}\) is instance~(i) of
\cref{def:p02-w02-p02-038-transport} for \(\Delta=\mathrm{Var}_{\mathrm{MS}}\),
and \(s^{\sharp}\) transports \(s\).  By clause~3 of
\cref{thm:p02-w02-p02-038-transport}, \(\mathcal M^{\sharp}\) satisfies
\((\mathrm{Ct})\) and \((\mathrm{Cl})\), including at tuples containing
\(\star\) or a wrongly tagged element, where the antecedent fails; it
satisfies \((\mathrm{Ne})\) because every carrier of the strict
\(\mathcal M\) is nonempty; and it satisfies \((\mathrm{Dj})\) because the
tagged images are pairwise disjoint.  By clause~4,
\(s^{\sharp}\) satisfies
\(\mathsf{Guard}(\mathrm{Var}_{\mathrm{MS}})=\mathsf{VarGuard}_{\mathrm{MS}}\).
By clause~2 of that lemma, for strict-language formulas,
\(\mathcal M,s\entails\varphi\) iff
\(\mathcal M^{\sharp},s^{\sharp}\entails\operatorname{tr}_{\mathrm{MS}}(\varphi)\).

\emph{Clause 3, left to right.}  Suppose
\(\Gamma\entails_{\mathrm{MS}}\varphi\), and let \(\mathcal A,\sigma\) satisfy
\(\mathsf{Ax}_{\mathrm{sort}}\), \(\mathsf{VarGuard}_{\mathrm{MS}}\) and
\(\operatorname{tr}_{\mathrm{MS}}[\Gamma]\).  By clause~1,
\(\mathcal A^{\downarrow},\sigma{\restriction}\mathrm{Var}_{\mathrm{MS}}\) is a
strict structure with a total assignment satisfying every member of
\(\Gamma\); hence it satisfies \(\varphi\), and by clause~1 again
\(\mathcal A,\sigma\entails\operatorname{tr}_{\mathrm{MS}}(\varphi)\).

\emph{Clause 3, right to left.}  Suppose the one-sorted consequence holds,
and let \(\mathcal M\) be strict and \(s\) total with \(\mathcal M,s\entails
\gamma\) for every \(\gamma\in\Gamma\).  By clause~2,
\(\mathcal M^{\sharp},s^{\sharp}\) satisfies \(\mathsf{Ax}_{\mathrm{sort}}\),
\(\mathsf{VarGuard}_{\mathrm{MS}}\) and
\(\operatorname{tr}_{\mathrm{MS}}[\Gamma]\), so it satisfies
\(\operatorname{tr}_{\mathrm{MS}}(\varphi)\), and by clause~2 again
\(\mathcal M,s\entails\varphi\).  No choice is used in either direction.
\end{proof}
\end{annexitem}

\begin{annexitem}{annex:p02-w02-p02-038-coverage-reduction}{thm:p02-w02-p02-038-coverage-reduction}{The coverage variant of the reduction}
\begin{proof}
\emph{Left to right.}  If \(\Gamma\entails_{\mathrm{MS}}\varphi\), then by
\cref{thm:p02-w02-p02-038-strict-correspondence}, clause~3, the translated
consequence holds from \(\mathsf{Ax}_{\mathrm{sort}}\cup
\mathsf{VarGuard}_{\mathrm{MS}}\cup\operatorname{tr}_{\mathrm{MS}}[\Gamma]\);
adding the premise \(\mathsf{Cov}\) keeps it, since every pair satisfying
the larger premise set satisfies the smaller one.

\emph{Right to left.}  Suppose the one-sorted consequence with
\(\mathsf{Cov}\) holds, and let \(\mathcal M\) be strict and \(s\) total with
\(\mathcal M,s\entails\gamma\) for every \(\gamma\in\Gamma\).  Write
\(\mathsf{Sort}=\{\mathsf s_1,\ldots,\mathsf s_k\}\) as in
\cref{def:p02-w02-p02-038-coverage}, and let the fixed element for one-sorted
variables outside the image of the variable injection be
\((\mathsf s_1,s(x^{\mathsf s_1}))\).  Form \(\mathcal M^{\sharp}_{s}\),
instance~(ii) of \cref{def:p02-w02-p02-038-transport}, with the assignment
\(\sigma\) of that instance, which transports \(s\).  By clause~3 of
\cref{thm:p02-w02-p02-038-transport} and the checks of instance~(ii),
\(\mathcal M^{\sharp}_{s}\) satisfies \((\mathrm{Ct})\), \((\mathrm{Cl})\),
\((\mathrm{Ne})\) (every carrier is nonempty), \((\mathrm{Dj})\) (tagged images
are disjoint) and \(\mathsf{Cov}\) (the images cover the domain).  By
clause~4, \(\sigma\) satisfies \(\mathsf{VarGuard}_{\mathrm{MS}}\), and by
clause~2, \(\mathcal M^{\sharp}_{s},\sigma\) satisfies
\(\operatorname{tr}_{\mathrm{MS}}[\Gamma]\).  So it satisfies
\(\operatorname{tr}_{\mathrm{MS}}(\varphi)\), and clause~2 returns
\(\mathcal M,s\entails\varphi\).

\emph{The remaining assertions.}  Both one-sorted consequence relations, with
and without \(\mathsf{Cov}\), are equivalent to
\(\Gamma\entails_{\mathrm{MS}}\varphi\), by this theorem and by
\cref{thm:p02-w02-p02-038-strict-correspondence}, clause~3; so they are
equivalent to each other.  The structure \(\mathcal M^{\sharp}\) does not
satisfy \(\mathsf{Cov}\), since its element \(\star\) lies in no
\(Q_{\mathsf s}^{\mathcal M^{\sharp}}\) (instance~(i)).
\end{proof}
\end{annexitem}

\begin{annexitem}{annex:p02-w02-p02-038-total-collapse}{thm:p02-w02-p02-038-total-collapse}{Total assignments hide empty sorts}
\begin{proof}
Let \(s\) be a total assignment into an E-structure \(\mathcal M\).  For each
sort \(\mathsf s\), \(s(x^{\mathsf s})\in M_{\mathsf s}\), so every carrier is
nonempty; the carriers of an E-structure are pairwise disjoint; so
\(\mathcal M\) is strict.  Every strict structure is an E-structure.  Hence
the pairs \((\mathcal M,s)\) of an E-structure and a total assignment are
exactly the pairs of a strict structure and a total assignment, and the two
consequence relations quantify over the same pairs with the same
satisfaction relation (\cref{def:p02-w02-p02-038-satisfaction}).
\end{proof}
\end{annexitem}

\begin{annexitem}{annex:p02-w02-p02-038-inhabited}{thm:p02-w02-p02-038-inhabited}{Sorts that can be empty}
\begin{proof}
Let \(I\) be the set of sorts inhabited by the signature.  If
\(\mathsf s\in I\), let \(t\) be a variable-free term of sort \(\mathsf s\);
in any E-structure \(\mathcal M\), its value under the empty assignment (the
\(\varnothing\)-assignment) is an element of \(M_{\mathsf s}\), so
\(M_{\mathsf s}\ne\varnothing\).  This proves that a sort which is empty in
some E-structure is not in \(I\).

For the converse and the final assertion, define \(\mathcal M_0\) by
\(M_{\mathsf s}\defeq\{\mathsf s\}\) for \(\mathsf s\in I\) and
\(M_{\mathsf s}\defeq\varnothing\) for \(\mathsf s\notin I\).  For
\(c:\longrightarrow\mathsf s\), the term \(c\) has no variable, so
\(\mathsf s\in I\); put \(c^{\mathcal M_0}\defeq\mathsf s\).  For
\(f:\mathsf s_1\times\cdots\times\mathsf s_n\longrightarrow\mathsf s\) with
\(n\geq1\): if some \(\mathsf s_i\notin I\), the product
\(M_{\mathsf s_1}\times\cdots\times M_{\mathsf s_n}\) is empty and
\(f^{\mathcal M_0}\) is the empty function; if every \(\mathsf s_i\in I\),
there are variable-free terms \(t_i\) of sort \(\mathsf s_i\), so
\(f(t_1,\ldots,t_n)\) is a variable-free term of sort \(\mathsf s\), hence
\(\mathsf s\in I\), and \(f^{\mathcal M_0}\) sends the only tuple
\((\mathsf s_1,\ldots,\mathsf s_n)\) to \(\mathsf s\).  Every relation
symbol is interpreted by the empty set.  The carriers are pairwise disjoint,
since \(\{\mathsf s\}\cap\{\mathsf t\}=\varnothing\) for
\(\mathsf s\ne\mathsf t\).  So \(\mathcal M_0\) is an E-structure, and
\(M_{\mathsf s}=\varnothing\) exactly when \(\mathsf s\notin I\).
\end{proof}
\end{annexitem}

\begin{annexitem}{annex:p02-w02-p02-038-empty-reduction}{thm:p02-w02-p02-038-empty-reduction}{Reduction for the empty-sort variant}
\begin{proof}
\emph{Left to right.}  Suppose \(\Gamma\entails^{\Delta}_{\mathrm E}\varphi\),
and let \(\mathcal A,\sigma\) satisfy \(\mathsf{Ax}^{\mathrm E}_{\mathrm{sort}}\),
\(\mathsf{Guard}(\Delta)\) and \(\operatorname{tr}_{\mathrm{MS}}[\Gamma]\).
By \cref{thm:p02-w02-p02-038-restriction}, clause~1, \(\mathcal A^{\downarrow}\)
is an E-structure (from \((\mathrm{Ct})\), \((\mathrm{Cl})\) and
\((\mathrm{Dj})\)) and \(g\defeq\sigma{\restriction}\Delta\) is a
\(\Delta\)-assignment into it; nonemptiness of the carriers is not needed.
By clause~3 of that lemma, \(\mathcal A^{\downarrow}\entails\gamma[g]\) for
every \(\gamma\in\Gamma\).  Hence \(\mathcal A^{\downarrow}\entails\varphi[g]\),
and clause~3 gives
\(\mathcal A,\sigma\entails\operatorname{tr}_{\mathrm{MS}}(\varphi)\).

\emph{Right to left.}  Suppose the one-sorted consequence holds, and let
\(\mathcal M\) be an E-structure and \(g\) a \(\Delta\)-assignment with
\(\mathcal M\entails\gamma[g]\) for every \(\gamma\in\Gamma\).  Form
\(\mathcal M^{\sharp}\), instance~(i) of \cref{def:p02-w02-p02-038-transport},
with the assignment \(\sigma\) sending \(v\in\Delta\) of sort \(\mathsf s\) to
\((\mathsf s,g(v))\) and every other one-sorted variable to \(\star\); it
transports \(g\).  The domain contains \(\star\), so it is nonempty even when
every carrier is empty.  By clause~3 of \cref{thm:p02-w02-p02-038-transport},
\(\mathcal M^{\sharp}\) satisfies \((\mathrm{Ct})\), \((\mathrm{Cl})\) and,
the tagged images being disjoint, \((\mathrm{Dj})\).  By clause~4,
\(\sigma\) satisfies \(\mathsf{Guard}(\Delta)\).  By clause~2 (strict
language), \(\mathcal M^{\sharp},\sigma\entails
\operatorname{tr}_{\mathrm{MS}}(\gamma)\) for every \(\gamma\in\Gamma\).  So
\(\mathcal M^{\sharp},\sigma\entails\operatorname{tr}_{\mathrm{MS}}(\varphi)\),
and clause~2 gives \(\mathcal M\entails\varphi[g]\).
\end{proof}
\end{annexitem}

\begin{annexitem}{annex:p02-w02-p02-038-empty-comparison}{thm:p02-w02-p02-038-empty-comparison}{Empty-sort consequence compared with the strict convention}
\begin{proof}
\emph{Clause 1.}  Let \(\mathcal M\) be strict and \(s\) total with
\(\mathcal M,s\entails\gamma\) for every \(\gamma\in\Gamma\).  \(\mathcal M\) is
an E-structure and \(s{\restriction}\Delta\) is a \(\Delta\)-assignment.  By
\cref{thm:p02-w02-p02-038-coincidence}, \(\mathcal M\entails
\gamma[s{\restriction}\Delta]\) for every \(\gamma\in\Gamma\), so
\(\mathcal M\entails\varphi[s{\restriction}\Delta]\), and coincidence again
gives \(\mathcal M,s\entails\varphi\).

\emph{Clause 2.}  By \cref{thm:p02-w02-p02-038-inhabited} there is an
E-structure \(\mathcal M_0\) with \(M_{\mathsf s}=\varnothing\).  Both
\(\forall v\,\psi\) and \(\exists v\,\psi\) have no free variable, so the
context \(\varnothing\) is admissible; let \(g_{\varnothing}\) be the empty
assignment.  There is no \(b\in M_{\mathsf s}\), so
\(\mathcal M_0\entails\forall v\,\psi[g_{\varnothing}]\) and
\(\mathcal M_0\not\entails\exists v\,\psi[g_{\varnothing}]\).  For the strict
semantics, let \(\mathcal M,s\entails\forall v\,\psi\); the carrier
\(M_{\mathsf s}\) is nonempty and contains \(b\defeq s(v)\), and
\(\mathcal M,s[v\mapsto b]\entails\psi\), so
\(\mathcal M,s\entails\exists v\,\psi\).  For \(\psi\defeq(v=v)\), the
premise \(\forall v\,(v=v)\) holds at every assignment in every sorted
structure, so the two results say that \(\exists v\,(v=v)\) fails at
\(\mathcal M_0,g_{\varnothing}\) while it holds at every strict pair.

\emph{Clause 3.}  Let \(\Delta\subseteq\Delta'\) and let \(g'\) be a
\(\Delta'\)-assignment into an E-structure with every member of \(\Gamma\)
true at \(g'\).  By coincidence, every member of \(\Gamma\) is true at
\(g'{\restriction}\Delta\); by the hypothesis, \(\varphi\) is true there; by
coincidence, \(\varphi\) is true at \(g'\).  For the failure of shrinking: if
\(g\) is a \(\{w\}\)-assignment into an E-structure, then
\(b\defeq g(w)\in M_{\mathsf s}\) and \(\mathcal M\entails(v=v)[g[v\mapsto b]]\),
so \(\exists v\,(v=v)\) is true at \(g\).  At the context \(\varnothing\),
clause~2 gives the countermodel \(\mathcal M_0,g_{\varnothing}\).

\emph{Clause 4.}  If every sort is inhabited by the signature, the first
paragraph of the proof of \cref{thm:p02-w02-p02-038-inhabited} makes every
carrier of every E-structure nonempty, so every E-structure is strict.  The
implication from \(\Gamma\entails^{\Delta}_{\mathrm E}\varphi\) to
\(\Gamma\entails_{\mathrm{MS}}\varphi\) is clause~1.  For the converse,
suppose \(\Gamma\entails_{\mathrm{MS}}\varphi\), and let \(g\) be a
\(\Delta\)-assignment into an E-structure \(\mathcal M\) with every member of
\(\Gamma\) true at \(g\).  The symbol set is at most countable
(\cref{def:p02-many-sorted-signature}); fix one injection of it into the
natural numbers.  It induces an injective coding of finite strings of
symbols and punctuation by natural numbers, and for each sort \(\mathsf s\)
let \(t_{\mathsf s}\) be the variable-free term of sort \(\mathsf s\) with
the least code.  Define the total assignment \(s\) by \(s(v)\defeq g(v)\) for
\(v\in\Delta\) and \(s(v)\defeq\llbracket t_{\mathsf s}\rrbracket^{\mathcal M}
_{g_{\varnothing}}\) for \(v\notin\Delta\) of sort \(\mathsf s\).  By
coincidence every member of \(\Gamma\) is true at \(s\); so \(\varphi\) is
true at \(s\), and by coincidence at \(g=s{\restriction}\Delta\).  No choice
principle is used beyond fixing the one injection.

\emph{Clause 5.}  Right to left: given strict \(\mathcal M\) and total \(s\)
with \(\Gamma\) true, the argument of clause~1 applies verbatim, and uses no
choice.  Left to right: suppose \(\Gamma\entails_{\mathrm{MS}}\varphi\), and
let \(g\) be a \(\Delta\)-assignment into a strict \(\mathcal M\) with
\(\Gamma\) true at \(g\).  Let \(F\) be a choice function on the family
\((M_{\mathsf s})_{\mathsf s\in\mathsf{Sort}}\) of nonempty sets; it exists
by the Axiom of Choice, and for finite \(\mathsf{Sort}\) by finitely many
existential instantiations.  Put \(s(v)\defeq g(v)\) for \(v\in\Delta\) and
\(s(v)\defeq F(\mathsf s)\) for \(v\notin\Delta\) of sort \(\mathsf s\).  By
coincidence \(\Gamma\) is true at \(s\), so \(\varphi\) is true at \(s\), and
by coincidence at \(g\).
\end{proof}
\end{annexitem}

\begin{annexitem}{annex:p02-w02-p02-038-overlap-reduction}{thm:p02-w02-p02-038-overlap-reduction}{One-sorted reduction for overlapping sorts}
\begin{proof}
\emph{Left to right.}  Suppose \(\Gamma\entails_{\cup}\varphi\), and let
\(\mathcal A,\sigma\) satisfy \(\mathsf{Ax}^{\cup}_{\mathrm{sort}}\),
\(\mathsf{VarGuard}_{\mathrm{MS}}\) and
\(\operatorname{tr}_{\mathrm{MS}}[\Gamma]\).  By
\cref{thm:p02-w02-p02-038-restriction}, clause~1, \(\mathcal A^{\downarrow}\)
is a \(\cup\)-structure (from \((\mathrm{Ne})\), \((\mathrm{Ct})\) and
\((\mathrm{Cl})\); disjointness is not used) and
\(g\defeq\sigma{\restriction}\mathrm{Var}_{\mathrm{MS}}\) is total.  Clause~3
of that lemma covers the cross-sort language, so every member of \(\Gamma\)
is true in \(\mathcal A^{\downarrow}\) at \(g\); hence \(\varphi\) is, and
clause~3 gives \(\mathcal A,\sigma\entails\operatorname{tr}_{\mathrm{MS}}(\varphi)\).

\emph{Right to left.}  Suppose the one-sorted consequence holds, and let
\(\mathcal M\) be a \(\cup\)-structure and \(s\) total with every member of
\(\Gamma\) true.  Form \(\mathcal M^{\cup}\), instance~(iii) of
\cref{def:p02-w02-p02-038-transport}, with \(\Delta=\mathrm{Var}_{\mathrm{MS}}\)
and the assignment \(\sigma\) sending each sorted \(v\) to \(s(v)\) and every
other one-sorted variable to \(U\); it transports \(s\).  By clause~3 of
\cref{thm:p02-w02-p02-038-transport}, \(\mathcal M^{\cup}\) satisfies
\((\mathrm{Ct})\), \((\mathrm{Cl})\) and \((\mathrm{Ne})\), every carrier being
nonempty.  By clause~4, \(\sigma\) satisfies
\(\mathsf{VarGuard}_{\mathrm{MS}}\).  The inclusion family is coherent, so
clause~2 applies to the cross-sort language and gives
\(\mathcal M^{\cup},\sigma\entails\operatorname{tr}_{\mathrm{MS}}(\gamma)\) for
every \(\gamma\in\Gamma\).  Hence
\(\mathcal M^{\cup},\sigma\entails\operatorname{tr}_{\mathrm{MS}}(\varphi)\), and
clause~2 gives \(\mathcal M,s\entails\varphi\).
\end{proof}
\end{annexitem}

\begin{annexitem}{annex:p02-w02-p02-038-overlap-comparison}{thm:p02-w02-p02-038-overlap-comparison}{Overlapping sorts compared with the strict convention}
\begin{proof}
\emph{Clause 1.}  Every strict structure is a \(\cup\)-structure, so
\(\Gamma\entails_{\cup}\varphi\) implies \(\Gamma\entails_{\mathrm{MS}}\varphi\).
Conversely, suppose \(\Gamma\entails_{\mathrm{MS}}\varphi\), and let
\(\mathcal M\) be a \(\cup\)-structure and \(s\) total with every member of
\(\Gamma\) true.  Form the tagged \(\mathcal M^{\sharp}\) (instance~(i), with
\(\Delta=\mathrm{Var}_{\mathrm{MS}}\)) and the transporting assignment
\(s^{\sharp}\).  By clause~3 of \cref{thm:p02-w02-p02-038-transport} it
satisfies \((\mathrm{Ct})\), \((\mathrm{Cl})\), \((\mathrm{Ne})\) (carriers
nonempty) and \((\mathrm{Dj})\) (tagged images disjoint, whether or not the
carriers overlap); by clause~4, \(s^{\sharp}\) satisfies
\(\mathsf{VarGuard}_{\mathrm{MS}}\); by clause~2 for the strict language,
which needs no coherence, \(s^{\sharp}\) satisfies every translated member
of \(\Gamma\).  By \cref{thm:p02-w02-p02-038-strict-correspondence},
clause~3, the translated consequence holds, so
\(\mathcal M^{\sharp},s^{\sharp}\entails\operatorname{tr}_{\mathrm{MS}}(\varphi)\),
and clause~2 gives \(\mathcal M,s\entails\varphi\).

\emph{Clause 2.}  Let \(\mathcal M\) have every carrier equal to \(\{0\}\),
every constant and every function value equal to \(0\), and every relation
empty.  It is a \(\cup\)-structure, and \(\theta\) holds in it with both
witnesses equal to \(0\).  The translation is
\[
 \operatorname{tr}_{\mathrm{MS}}(\theta)=
 \exists x^{\mathsf s}\bigl(Q_{\mathsf s}(x^{\mathsf s})\land
 \exists x^{\mathsf t}(Q_{\mathsf t}(x^{\mathsf t})\land
 x^{\mathsf s}=x^{\mathsf t})\bigr).
\]
In a model of \((\mathrm{Dj})\), witnesses would be one element lying in
\(Q_{\mathsf s}\) and in \(Q_{\mathsf t}\), which the disjointness instance
for \(\mathsf s,\mathsf t\) excludes; so
\(\mathsf{Ax}_{\mathrm{sort}}\entails_{\mathrm{FOL}}
\neg\operatorname{tr}_{\mathrm{MS}}(\theta)\), and
\(\operatorname{tr}_{\mathrm{MS}}(\neg\theta)=
\neg\operatorname{tr}_{\mathrm{MS}}(\theta)\).  If \((\mathrm{Dj})\) were
added to the premises of \cref{thm:p02-w02-p02-038-overlap-reduction}, the
resulting equivalence for \(\Gamma=\varnothing\) and the conclusion
\(\neg\theta\) would have a true right side, and so would assert
\(\varnothing\entails_{\cup}\neg\theta\), which \(\mathcal M\) refutes.

\emph{Clause 3.}  For the same \(\mathcal M\), the tagged structure has
\(Q_{\mathsf s}^{\mathcal M^{\sharp}}=\{(\mathsf s,0)\}\) and
\(Q_{\mathsf t}^{\mathcal M^{\sharp}}=\{(\mathsf t,0)\}\).  The translation of
\(\theta\) requires one element in both, but \((\mathsf s,0)\ne(\mathsf t,0)\)
because \(\mathsf s\ne\mathsf t\).  So
\(\mathcal M^{\sharp}\not\entails\operatorname{tr}_{\mathrm{MS}}(\theta)\)
although \(\mathcal M\entails\theta\); the tagged family is not coherent here,
as instance~(i) predicts.
\end{proof}
\end{annexitem}
